\chapter{Navier--Stokes, and what a formalization checks}

\textbf{Purpose:} Record what the official problem asks, what was published against it on
2026-09-08, and the result of checking the published formalization's statement against the problem
statement point by point. \textbf{Scope:} Fefferman's Clay problem description, 6 pages; OpenAI,
\emph{Finite time blowup for Navier--Stokes}, 166 pages; the Lean~4 repository
\texttt{openai/\allowbreak{}NavierStokesAndEuler}, four source files of 2{,}669 read.

\section{What the problem asks}

Fefferman's official description sets up the incompressible equations on $\R^n$ for $n=2$ or $3$,
with viscosity $\nu > 0$, and imposes decay and regularity conditions numbered (4) through (11). It
then states, in the sentence that governs everything after it:

\begin{quote}
To give reasonable leeway to solvers while retaining the heart of the problem, we ask for a proof
of one of the following four statements.
\end{quote}

The four are (A) existence and smoothness on $\R^3$, (B) the same on $\R^3/\Z^3$, (C) breakdown on
$\R^3$, and (D) breakdown on $\R^3/\Z^3$. In (A) and (B) the force $f$ is taken identically zero. In
(C) and (D) it is not. Statement (C) reads, in his words, that for $\nu>0$ and $n=3$ there exist a
smooth divergence-free $u^\circ$ on $\R^3$ \emph{and a smooth} $f(x,t)$ on
$\R^3\times[0,\infty)$, satisfying (4) and (5), for which there exist no solutions $(p,u)$ of (1),
(2), (3), (6), (7) on $\R^3\times[0,\infty)$. Statement (D) is the same with (8) and (9) on the data
and (10) and (11) on the solution.

So a smooth external force is named inside two of the four statements, and a proof of any one of
the four is what the problem asks for. The only occurrence of the word \emph{prize} in the document
is a remark that the Euler equations are not on the Clay Institute's list. There is no passage
reserving the problem for (A) or (B).

The conditions that matter for the check below are (4), rapid decay of every derivative of the
initial data; (5), rapid decay of every space-time derivative of the force on the closed half-line;
(6), smoothness of $p$ and $u$ on $\R^n\times[0,\infty)$; and (7), $\int|u(x,t)|^2\,dx$ bounded by
one constant for all $t\ge0$, which Fefferman annotates \emph{bounded energy}.

\textbf{The document carries errata and they are load-bearing.} The first reads that the further
condition $p(x+e_j,t)=p(x,t)$ should be made explicit in Eqn (8). The periodic problem therefore
requires a periodic pressure, which the body of the statement does not say.

\section{What was published against it}

On 2026-09-08 OpenAI published \emph{Finite time blowup for Navier--Stokes}, 166 pages, together
with a companion Euler paper and a Lean~4 repository. Its Theorem 1.1 is a construction, not
directly a non-existence statement:

\begin{quote}
For every $\nu>0$ there exist a force $f \in C_c^\infty(\R^3\times(0,1);\R^3)$, a compact set
$K\subset\R^3$, and smooth velocity and pressure fields $u,p$ on $\R^3\times[0,1)$ satisfying
$\partial_t u + (u\cdot\nabla)u - \nu\Delta u + \nabla p = f$, $\nabla\cdot u = 0$,
$u(\cdot,0)=0$, such that $\operatorname{supp} u(\cdot,t)\cup\operatorname{supp}p(\cdot,t)\subseteq
K$ for every $0\le t<1$, $\sup_t \|u(t)\|_{L^2(\R^3)} < \infty$, and
$\limsup_{t\to1}\|u(t)\|_{L^\infty(\R^3)} = \infty$.
\end{quote}

Constructing one solution that blows up does not by itself establish (C). Statement (C) asks that
\emph{no} smooth bounded-energy solution exist on the whole half-line for that datum and force, and
a blowing-up solution leaves open whether some other smooth global one exists. The paper closes
that gap and says so in the sentence after the theorem, asserting that consequently there is no
smooth solution on $\R^3\times[0,\infty)$ with the same force and initial datum whose kinetic
energy is uniformly bounded, and that this establishes alternative (C) as stated by Fefferman.

The argument for the gap is in the proof of Theorem 1.1 in section 10.4. Smoothness and fixed
compact support put the field in $C([0,T];H^3)$ for each $T<1$; the embedding
$H^3(\R^3)\hookrightarrow L^\infty(\R^3)$ excludes a classical $H^3$ continuation through time one;
a difference-energy calculation with Gr\"onwall gives uniqueness on every shorter interval, making
$[0,1)$ the maximal classical existence interval. A supposed global smooth competitor of uniformly
bounded energy is then forced by the comparison lemma to agree with the constructed field
throughout $[0,1)$, which contradicts its own boundedness. Every fixed positive viscosity is
reached by a spatial rescaling that leaves the singular time at one.

Statement (D) is Corollary 10.6, obtained by rescaling the compact supports into the interior of a
fundamental cube and summing integer translates, whose supports stay disjoint so the nonlinear term
survives. The corollary's proof states that the pressure is periodic as well, \emph{as required in
the erratum to the problem statement}. The authors read the errata.

\section{What a proof assistant does not check}

This project's position on machine-checked proof, stated before this chapter was written, is that
Lean verifies the deduction and never the claim. A proof assistant confirms that a conclusion
follows from the assumptions as written. It cannot confirm that the statement formalized is the
statement anybody meant. A machine-checked proof of a neighbouring theorem compiles exactly as
cleanly as a machine-checked proof of the right one, and nothing in the compile distinguishes them.
The burden therefore moves entirely onto the formalized statement, and that statement is a
translation, which is the part no proof assistant checks.

That makes the formalized statement a finite document a person reads word by word. The rest of this
chapter is that reading.

\section{The statement check}

Six points were checked against the Lean statement in
\texttt{ComparatorChallenges/\allowbreak{}NavierStokes.\allowbreak{}lean} and against Fefferman's
numbered conditions. The first five were specified before the file was opened. The sixth was added
because the file itself raised it.

\begin{center}
\begin{tabular}{@{}p{0.30\textwidth}p{0.44\textwidth}c@{}}
\toprule
Point & What the Lean statement carries & Verdict \\
\midrule
Viscosity strictly positive, and not a vanishing-viscosity limit
& \texttt{(nu : R) (hnu : nu > 0)}, universally quantified over the theorem
& matches \\[0.5em]
Condition (4) on the initial data, unrelaxed
& \texttt{InitialVelocityConditionDecay}: divergence-free, $C^\infty$, and for every derivative
order $m$ and every $K$ a constant $C$ with
$\|\nabla^m u_0(x)\| \le C/(1+\|x\|)^K$
& matches \\[0.5em]
Condition (5) on the force, unrelaxed, on the closed half-line
& \texttt{ForceConditionDecay}: $C^\infty$ on $\mathrm{univ}\times[0,\infty)$ and
$\|\nabla^m f(x,t)\| \le C/(1+\|x\|+t)^K$ for all $t\ge0$
& matches \\[0.5em]
Conclusion in the negated existential form Fefferman quantifies
& \texttt{not (exists v p, NavierStokesExistence\-AndSmoothnessRn nu u0 f v p)}
& matches \\[0.5em]
Energy bounded uniformly in time, not at $t=0$ alone
& \texttt{globally\_bounded\_energy}: there exists $E$ such that for all $t\ge0$,
$\int\|v(x,t)\|^2 < E$; plus square-integrability at each $t\ge0$
& matches \\[0.5em]
The periodic pressure condition from the errata
& \texttt{isOnePeriodic\_pressure} on the periodic structure, annotated in the source as following
the Clay errata
& matches \\
\bottomrule
\end{tabular}
\end{center}

The equation itself is encoded with the time derivative taken as a derivative within
$[0,\infty)$, which is the right treatment given that Fefferman states equation (1) on the closed
half-line. Divergence is defined as the trace of the Jacobian and carries an explicit junk value of
zero at points of non-differentiability, which is declared in the source rather than left implicit.

\section{The statement is not the authors' own, and that is the strongest fact here}

The file carrying the statement is a \emph{comparator challenge}. Its header records that it was
copied, at a named commit, from the Formal Conjectures project's Lean formalization of the
Navier--Stokes problem statement, that the breakdown alternatives (C) and (D) are retained
\emph{including their intentional placeholder gaps}, and that neither the proof root nor the
submission imports it. Both theorems in that file end in a placeholder and prove nothing. It is a
yardstick and not a result.

The discharge is a separate module. It restates both theorems with textually identical statements
and closes each by applying a bridge lemma from the project's own development, declaring no axioms
and leaving no placeholder, and it ends by printing the axiom dependencies of both theorems. The
repository's instructions direct a checker to verify the two challenges with an external checking
tool under a sandbox, using an independent exporter and an independent kernel rather than Lean's
own.

That architecture is a direct answer to the objection this chapter opened with. The concern is that
an author who writes both the statement and the proof can machine-check a translation of their own
choosing. Here the statement was written by a third party, published beforehand, for the purpose of
stating the problem and not of being proved; the authors adopted it unmodified and discharged it;
and the axiom base is printed rather than asserted. The translation is still the thing to read, and
this chapter read it, but the party that wrote it had no stake in what would later be proved
against it.

\section{What was not verified, stated here and not in a footnote}

\textbf{Nothing here was machine-checked.} The Lean development was not built. No proof was
executed, no axiom output was observed, and the claim that the discharge module compiles without
placeholders rests on reading its source text rather than on running it. The check in this chapter
is a reading of four files out of 2{,}669.

\textbf{The proof was not read.} The 166-page paper was read at its abstract, its introduction, its
Theorem 1.1, the section 10.4 argument that bridges to the non-existence form, and Corollary 10.6.
The mathematical content of sections 3 through 9 and the three appendices was not read. The bridge
lemma the discharge module applies was not traced into the development.

\textbf{The priority dispute was not investigated.} A dispute over credit is running in this field.
It was not examined for this book, no position is taken on it, and nothing in this chapter bears on
it either way.

\textbf{One reported fact is recorded as reported.} That the authors have stated they will not claim
the Millennium Prize appeared in a headline encountered during the search and was not confirmed
against any statement of theirs. It is written here as an unverified report and should not be
repeated from this book as anything else.

\textbf{And the status question is separate from the mathematics question.} Whether the Clay
Mathematics Institute has accepted, awarded, or will award anything was not verified and is not
asserted. What this chapter establishes is narrower and is the part that is checkable: the
statement discharged in Lean is the statement Fefferman wrote, on all six points examined.
